% ECONOMICS_PROBLEM: {"id": "C31-129", "title": "Causal inference under interference and network spillovers", "jel_code": "C31", "source_id": "OP-0101"}
% FIRST_POSTED: 2026-10-09
% ATTEMPT_STATUS: unresolved
% ENTRY_KIND: research_attempt
\documentclass[11pt,a4paper]{article}
\usepackage[T1]{fontenc}
\usepackage[utf8]{inputenc}
\usepackage{lmodern,amsmath,amssymb,amsthm,mathtools}
\usepackage[margin=25mm]{geometry}
\usepackage{hyperref}
\hypersetup{colorlinks=true,urlcolor=blue,linkcolor=blue}
\setlength{\parindent}{0pt}
\setlength{\parskip}{5pt}
\newtheorem{theorem}{Theorem}
\newtheorem{proposition}[theorem]{Proposition}
\newtheorem{lemma}[theorem]{Lemma}
\newtheorem{corollary}[theorem]{Corollary}
\newcommand{\R}{\mathbb R}
\newcommand{\E}{\mathbb E}
\newcommand{\Prb}{\mathbb P}
\newcommand{\ind}{\mathbf 1}
\newcommand{\TV}{\operatorname{TV}}
\newcommand{\KL}{\operatorname{KL}}
\newcommand{\Var}{\operatorname{Var}}
\newcommand{\OPT}{\operatorname{OPT}}
\newcommand{\diam}{\operatorname{diam}}
\newcommand{\supp}{\operatorname{supp}}
\DeclareMathOperator*{\argmax}{arg\,max}
\DeclareMathOperator*{\argmin}{arg\,min}
\begin{document}
\noindent\textbf{Problem ID:} \texttt{C31-129}\quad\textbf{Model:} GPT 6 Astra Ultra\quad\textbf{Version:} 1\par
\section*{Causal inference under interference and network spillovers}

\textbf{Status.} Unresolved for unrestricted learned interference. We give simultaneous honest inference over a finite local exposure class, an explicit dependence/complexity rate, and a nonvanishing-width example when dependence is unrestricted.

\subsection*{Short target and sampling assumptions}
Fix $n$ units, a pretreatment network and covariates. All probabilities below condition on these fixed objects. Assignments $Z_1,\ldots,Z_n$ are independent Bernoulli with known probabilities. Potential outcomes are fixed functions of assignments. For each unit specify a set $B_i\subseteq[n]$ such that $Y_i(z)$ depends only on $z_{B_i}$ and $|Y_i(z)|\leq B$ for every assignment. The locality restriction is imposed on the actual outcome, not only on its exposure approximation.

Let $\mathcal H$ be a prespecified finite set of $H$ exposure maps, each $h_i(z)$ depending only on $z_{B_i}$. For each $h$ fix two distinct exposure labels $e_h,e'_h$, and assume their probabilities $p_i^h(e_h),p_i^h(e'_h)$ are at least $p>0$. Set
\[
 \tau_h=\frac1n\sum_i\left\{\E[Y_i(Z)\mid h_i(Z)=e_h]
             -\E[Y_i(Z)\mid h_i(Z)=e'_h]\right\}.
\]
This is a design-specific contrast even for insufficient maps. A learner may select $\widehat h$ using the entire realized experiment, provided it selects from $\mathcal H$.

Join distinct units $i,j$ in a dependency graph when $B_i\cap B_j\ne\varnothing$. Let $\chi$ be any number of colors in a proper coloring of this graph; $\chi\leq\Delta+1$ for maximum graph degree $\Delta$. Define
\[
 \widehat\tau_h=\frac1n\sum_iY_i(Z)
 \left\{\frac{\ind\{h_i(Z)=e_h\}}{p_i^h(e_h)}
       -\frac{\ind\{h_i(Z)=e'_h\}}{p_i^h(e'_h)}\right\}.
\]

\begin{lemma}[Colored concentration]
Suppose random variables $X_1,\ldots,X_n$ are independent within each class of a partition into $\chi$ classes and each lies in an interval of length at most $L$. Then, for every $t>0$,
\[
 \Prb\left(\left|n^{-1}\sum_i(X_i-\E X_i)\right|>t\right)
 \leq2\exp\left(-\frac{2nt^2}{\chi L^2}\right).
\]
\end{lemma}
\begin{proof}
Write $S_c=\sum_{i\in c}(X_i-\E X_i)$. H\"older's inequality with exponents all equal to $\chi$ and the independent-variable Hoeffding moment bound give
\[
 \E e^{\lambda\sum_cS_c}
 \leq\prod_c(\E e^{\chi\lambda S_c})^{1/\chi}
 \leq\exp(\chi\lambda^2 nL^2/8).
\]
Markov's inequality with $\lambda=4t/(\chi L^2)$ gives the upper tail $e^{-2nt^2/(\chi L^2)}$. Apply the same argument to negatives and add the two tail probabilities. Independence between different color classes is not required.
\end{proof}

\begin{theorem}[A valid interval after exposure-map selection]
Under the assumptions above, with probability at least $1-\alpha$, every $h\in\mathcal H$ satisfies
\[
 |\widehat\tau_h-\tau_h|\leq
 r_n:=\frac{B}{p}\sqrt{\frac{2\chi\log(2H/\alpha)}n}.
\]
Consequently $[\widehat\tau_{\widehat h}-r_n,\widehat\tau_{\widehat h}+r_n]$ covers the selected contrast $\tau_{\widehat h}$ with the same probability, even when selection uses outcomes.
\end{theorem}
\begin{proof}
Each summand is between $-B/p$ and $B/p$, because the two exposure indicators are disjoint. Its expectation is the corresponding difference of conditional means by the definition of conditional expectation. Within a color class the $B_i$ are disjoint; the summands are therefore independent functions of disjoint collections of independently assigned treatments. Apply the lemma with $L=2B/p$, then take a union bound over $H$ maps. The simultaneous event implies the selected-map statement for every possible selection rule.
\end{proof}

\begin{corollary}[Stable responses with bounded exposure error]
For every $h$ for which functions $y_i^h$ satisfy
$|Y_i(z)-y_i^h(h_i(z))|\leq\varepsilon_h$ for all assignments, put
$\tau_h^*=n^{-1}\sum_i[y_i^h(e_h)-y_i^h(e'_h)]$.
On the simultaneous event of the theorem,
\[
 |\widehat\tau_h-\tau_h^*|\leq r_n+2\varepsilon_h
\]
for all such maps. In particular this covers an outcome-selected map if its error bound holds uniformly and is known.
\end{corollary}
\begin{proof}
Conditional averaging preserves the error bound at each exposure, so $|\tau_h-\tau_h^*|\leq2\varepsilon_h$. Add the simultaneous estimation bound. The error assumption is substantive and cannot be verified merely by selecting a map with a small observed residual.
\end{proof}

For an undirected network of maximum degree $d$, if each $B_i$ is its closed neighborhood, each has at most $d+1$ coordinates and any coordinate occurs in at most $d+1$ such sets. Hence each dependency-graph vertex has degree at most $(d+1)^2-1$. One obtains informative diameter whenever
\[
 \frac{(d+1)^2\log(2H/\alpha)}{np^2}\longrightarrow0,
 \qquad\sup_h\varepsilon_h\longrightarrow0
\]
for the stable-response target. This sufficient condition need not be necessary, and better variance-based bounds can improve the $p$ dependence.

\begin{proposition}[Overlap alone cannot make confidence sets shrink]
For every $n$ there is a local-outcome experiment with two exposures of probability $1/2$, bounded outcomes in $[0,1]$, and parameter $\theta\in\{0,1\}$, such that every confidence set $C$ satisfying $\Prb_\theta(\theta\in C)\geq1-\alpha$ for both parameters obeys
\[
 \E_0\diam C\geq\frac12-2\alpha.
\]
\end{proposition}
\begin{proof}
Use $B_i=\{1\}$ for all $i$, $Z_1\sim\mathrm{Ber}(1/2)$, and $Y_i(z)=\theta z_1$. Other assignments can be independent Bernoulli and are irrelevant. The exposure map is $h_i(z)=z_1$, and the average exposure contrast is exactly $\theta$. On event $E=\{Z_1=0\}$ the full data laws coincide under $\theta=0$ and $\theta=1$; both give $E$ probability $1/2$. Thus, under $P_0$, the probability of $E$ and failure to include 1 is at most $\alpha$, by transferring that identical subprobability law from $P_1$. The probability of failure to include 0 is also at most $\alpha$. With probability at least $1/2-2\alpha$, both points belong to $C$, forcing diameter at least one. The dependency graph is complete despite exposure overlap being fixed.
\end{proof}

\subsection*{Remaining gap and reference}
The result pays for outcome-based map selection through $\log H$ and for local dependence through $\chi$. It does not handle an unrestricted map learner, unknown global interference, treatment-induced network formation, or unidentified observational treatment assignment. The target after map selection is explicitly $\tau_{\widehat h}$; selecting a convenient map does not establish that its contrast equals a policy effect.

Peter M. Aronow and Cyrus Samii (2017), \emph{Estimating Average Causal Effects Under General Interference, with Application to a Social Network Experiment}, Annals of Applied Statistics 11(4), 1912--1947. \url{https://doi.org/10.1214/16-AOAS1005}; author preprint record \url{https://arxiv.org/abs/1305.6156}. The primary abstract and bibliographic record were verified: they establish the role of known exposure mappings and local dependence. The elementary finite-class concentration and lower-bound calculations are proved here and not attributed as their exact theorems.

\end{document}
